\documentclass[12pt,a4paper]{article}
\usepackage{ugmath}
\usepackage{float}
\usepackage{placeins}
\usepackage[labelfont=bf,labelsep=space]{caption}
\newcommand{\paren}[1]{\left( #1 \right)}
\newcommand{\alumno}{Ricardo León Martínez}
\newcommand{\materia}{Fisica I}
\newcommand{\profesor}{Lucero Uscanga Aguilera}
\newcommand{\tarea}{Tarea 1}
\newcommand{\fecha}{16/2/2026}

\begin{document}
    \begin{center}
    {\large\textbf{UNIVERSIDAD DE GUANAJUATO}}\\[0.3cm]
    {\normalsize\textbf{DIVISIÓN DE CIENCIAS NATURALES Y EXACTAS}}\\
    {\normalsize\textbf{CAMPUS GUANAJUATO}}\\[1cm]

    {\Large\textbf{\tarea\ (\materia)}}\\[1cm]
    \end{center}

    \textbf{Nombre:} \alumno \hfill 
    \textbf{Fecha:} \fecha \hfill 
    \textbf{Calificación:} \rule{3cm}{0.4pt} \\[0.3cm]

    \begin{mdframed}[style=mdbluebox,frametitle={Ejercicio 1}]
        La viga $AB$ es uniforme y tiene una masa de $100kg$. Descansa en sus extremos $A$ y $B$
        y soporta las masas como se indica en la figura. Calcular la reacción en los soportes.
    \end{mdframed}

    \begin{solucion}
        Sea $g$ la aceleración de la gravedad. Las fuerzas verticales que actúan sobre la viga son:
        \begin{itemize}
            \item Peso de la viga: $100g$, aplicado en su punto medio, a $3m$ de $A$.
            \item Peso de la carga de $300kg$: $300g$, a $2m$ de $A$.
            \item Peso de la carga de \( 500\text{ kg} \): \( 500g \), a \( 4\text{ m} \) de \( A \).
            \item Reacción en \( A \): \( R_A \), hacia arriba.
            \item Reacción en \( B \): \( R_B \), hacia arriba.
        \end{itemize}
        Del equilibrio de fuerzas verticales se tiene
        \[
            R_A + R_B = 100g + 300g + 500g = 900g. \tag{1}
        \]
        Tomando momentos respecto del punto \( A \) (positivo el sentido antihorario), la suma de momentos debe ser nula:
        \[
            - (300g)(2) - (100g)(3) - (500g)(4) + R_B(6) = 0.
        \]
        Esto es
        \[
            -600g - 300g - 2000g + 6R_B = 0,
        \]
        \[
            -2900g + 6R_B = 0,
        \]
        de donde
        \[
            6R_B = 2900g \quad \Rightarrow \quad R_B = \frac{2900g}{6} = \frac{1450g}{3}. \tag{2}
        \]
        Sustituyendo (2) en (1):
        \[
            R_A + \frac{1450g}{3} = 900g,
        \]
        \[
            R_A = 900g - \frac{1450g}{3} = \frac{2700g - 1450g}{3} = \frac{1250g}{3}. \tag{3}
        \]
        Tomando \( g = 9.81\ \text{m/s}^2 \), se obtiene
        \[
            R_A = \frac{1250 \times 9.81}{3} = 4087.5\ \text{N},
            \qquad
            R_B = \frac{1450 \times 9.81}{3} = 4741.5\ \text{N}.
        \]
        En kilogramos–fuerza (\( 1\ \text{kgf} = 9.81\ \text{N} \)):
        \[
            R_A \approx 416.7\ \text{kgf}, \qquad R_B \approx 483.3\ \text{kgf}.
        \]
        Por tanto, las reacciones en los apoyos son
        \[
            R_A = 4087.5\ \text{N},\quad R_B = 4741.5\ \text{N}.
        \]
    \end{solucion}

    \begin{mdframed}[style=mdbluebox,frametitle={Ejercicio 2}]
        La viga uniforme $AB$ de la figura tiene $4cm$ de largo y pesa $100kgf$. La viga puede
        rotar alrededor del punto fijo $C$. La viga reposa en el punto $A$. Un hombre que pesa
        $75kgf$ camina a lo largo de la viga, partiendo de $A$. Calcular la máxima distancia que
        el hombre puede caminar a partir de $A$ manteniendo el equilibrio. Representar
        la fuerza de reaccion en $A$ como una función de la distancia $x$.
    \end{mdframed}
\end{document}